% Options for packages loaded elsewhere
\PassOptionsToPackage{unicode}{hyperref}
\PassOptionsToPackage{hyphens}{url}
\PassOptionsToPackage{dvipsnames,svgnames,x11names}{xcolor}
%
\documentclass[
  11pt,
]{article}
\usepackage{amsmath,amssymb}
\usepackage{iftex}
\ifPDFTeX
  \usepackage[T1]{fontenc}
  \usepackage[utf8]{inputenc}
  \usepackage{textcomp} % provide euro and other symbols
\else % if luatex or xetex
  \usepackage{unicode-math} % this also loads fontspec
  \defaultfontfeatures{Scale=MatchLowercase}
  \defaultfontfeatures[\rmfamily]{Ligatures=TeX,Scale=1}
\fi
\usepackage{lmodern}
\ifPDFTeX\else
  % xetex/luatex font selection
  \setmainfont[]{DejaVu Serif}
  \setmonofont[]{DejaVu Sans Mono}
  \setmathfont[]{DejaVu Math TeX Gyre}
\fi
% Use upquote if available, for straight quotes in verbatim environments
\IfFileExists{upquote.sty}{\usepackage{upquote}}{}
\IfFileExists{microtype.sty}{% use microtype if available
  \usepackage[]{microtype}
  \UseMicrotypeSet[protrusion]{basicmath} % disable protrusion for tt fonts
}{}
\makeatletter
\@ifundefined{KOMAClassName}{% if non-KOMA class
  \IfFileExists{parskip.sty}{%
    \usepackage{parskip}
  }{% else
    \setlength{\parindent}{0pt}
    \setlength{\parskip}{6pt plus 2pt minus 1pt}}
}{% if KOMA class
  \KOMAoptions{parskip=half}}
\makeatother
\usepackage{xcolor}
\usepackage[margin=25mm]{geometry}
\usepackage{longtable,booktabs,array}
\usepackage{calc} % for calculating minipage widths
% Correct order of tables after \paragraph or \subparagraph
\usepackage{etoolbox}
\makeatletter
\patchcmd\longtable{\par}{\if@noskipsec\mbox{}\fi\par}{}{}
\makeatother
% Allow footnotes in longtable head/foot
\IfFileExists{footnotehyper.sty}{\usepackage{footnotehyper}}{\usepackage{footnote}}
\makesavenoteenv{longtable}
\setlength{\emergencystretch}{3em} % prevent overfull lines
\providecommand{\tightlist}{%
  \setlength{\itemsep}{0pt}\setlength{\parskip}{0pt}}
\setcounter{secnumdepth}{-\maxdimen} % remove section numbering
\setlength{\emergencystretch}{4em}
\sloppy
\setlength{\parskip}{4pt}
\ifLuaTeX
  \usepackage{selnolig}  % disable illegal ligatures
\fi
\usepackage{bookmark}
\IfFileExists{xurl.sty}{\usepackage{xurl}}{} % add URL line breaks if available
\urlstyle{same}
\hypersetup{
  pdftitle={A sharp uniform diameter threshold for the maximum Wiener index in a fixed-edge middle range},
  pdfauthor={Ryutaro Yonezu},
  colorlinks=true,
  linkcolor={black},
  filecolor={Maroon},
  citecolor={Blue},
  urlcolor={blue},
  pdfcreator={LaTeX via pandoc}}

\title{A sharp uniform diameter threshold for the maximum Wiener index
in a fixed-edge middle range}
\author{Ryutaro Yonezu}
\date{Research No.~89 --- v1.0 preprint, 28 November 2026}

\begin{document}
\maketitle
\begin{abstract}
For connected simple graphs with fixed order, diameter and edge count,
we establish an exact formula and unique extremal construction in a
strict middle surplus range above an explicit diameter threshold. The
threshold for the simultaneous value and uniqueness assertion is sharp.
The proof combines general localization and deficit estimates with
finite integer and exact branch-and-bound certificates. Independent AI
audit: CONDITIONAL PASS; not human peer reviewed. Novelty inconclusive;
references partial.
\end{abstract}

\section{1. Statement and scope}\label{statement-and-scope}

The Wiener index of a connected finite simple undirected graph \(G\) is
\(W(G)=\sum_{\{u,v\}\subset V(G)}\operatorname{dist}_G(u,v)\). We fix
the order \(n\), diameter \(d\) and edge count \(m\). Write

\[k=n-d-1,\quad r=m-n+1,\quad A=\binom{k+1}{2},\quad h=A-r,\quad M=\left\lfloor k^2/8\right\rfloor.\]

Throughout the main theorem, \(k\ge4\) and \(0<h<\lfloor k^2/4\rfloor\).
Set

\[R_0=\min\{j\in\{0,\ldots,\lfloor k/2\rfloor\}:j(k-j)\ge h\},\]
\[L_0=k-R_0,\quad q=L_0R_0-h,\quad F=\binom{d+2}{3}+k\binom{d+1}{2}.\]

\textbf{Theorem.} For every \(d\ge M+3\), the maximum Wiener index among
the graphs with these parameters is

\[\boxed{W_{\max}=F+hd+r-qR_0.}\]

The maximizing graph \(H_*\) is unique up to isomorphism. It consists of
an induced diametral path \(p_0,\ldots,p_d\), two off-path cliques of
sizes \(L_0,R_0\), all left vertices adjacent to \(p_0,p_1\), all right
vertices to \(p_{d-1},p_d\), and exactly \(q\) left vertices
additionally adjacent to \(p_2\). There are no other edges. The least
uniform diameter threshold for the value and uniqueness together over
the strict middle range is \(\lfloor k^2/8\rfloor+3\).

The construction is defined here for \(d\ge4\); the theorem uses
\(d\ge5\). Minimality of \(R_0\) gives \(0\le q\le L_0-R_0\) (Section
3). Its edge count is \(d+2k+q+\binom{L_0}{2}+\binom{R_0}{2}=d+k+r\).
The path is geodesic because all same-side excursions between path
vertices have length at least their coordinate difference, and the sides
have no cross edges. All other distances are at most \(d\), by routing
through \(p_1\) or \(p_{d-1}\) for ordinary vertices and \(p_2\) for
shortened vertices. Thus its diameter is \(d\). The identity below gives
its index.

The proof uses three disjoint regimes: a general analytic argument for
\(k\ge19\), a finite integer relaxation for \(11\le k\le18\), and an
exhaustive computer-assisted certificate for \(4\le k\le10\). No claim
is made for other surplus ranges or for a value-only threshold. In the
proof, \(H_{d,r}\) denotes the candidate graph (abbreviated to \(H_*\)),
while \(H\) without a subscript denotes the number of missing same-side
edges. We use \(\operatorname{binom}(a,b)=\binom{a}{b}\) and
\(\operatorname{floor}(x)=\lfloor x\rfloor\) in some inline expressions.
All path coordinates in a right group are reflected.

\section{2. Localization with touching
blocks}\label{localization-with-touching-blocks}

Fix a diametral geodesic P in a maximizing graph. Each connected
component C off P has \(s_{C}\) vertices and a nonempty attachment set
on P, with extreme coordinates \(\alpha _{C},\beta _{C}\). Its interval
width is at most \(s_{C}+1\) \(\le\) \(2s_{C}\): a route joining the
extreme attachments through C has length at most \(s_{C}+1\). Merge
intersecting closed intervals into blocks. A block B has s off-path
vertices, width \(w\le 2s\), and contains the full intervening P
segment. Its metric on the P segment is the path metric. Write local
coordinates 0..w.

For later use, the diameter of B is at most 2s. To see this, any vertex
of C can reach either extreme attachment in at most \(s_{C}\) steps. For
an off-path vertex and a path vertex inside its component interval,
choose the closer extreme, at distance at most
\(\operatorname{floor}(w_{C}/2)\le s_{C}\) along P. Outside that
interval the intervening gap is covered by the intervals of the
remaining components, whose total width is at most \(2(s-s_{C})\). For
vertices in different components with disjoint intervals, the gap is at
most \(2(s-s_{C}-s_{D})\). If their intervals intersect, some endpoint
of one lies in the other, and the distance to an endpoint of the other
is at most half the width of that other interval. More explicitly, if an
endpoint of C lies in the interval of D, the resulting route has length
at most \(s_{C}+\operatorname{floor}(w_{D}/2)+s_{D}\) \(\le\)
\(2(s_{C}+s_{D})\); reverse C and D if necessary. In both cases the
route length is \(\le 2s\). Vertices in the same off-path component have
distance \(\le s_{C}-1\). Path-path distances are \(\le w\).

Define \(a_B=\max_v\operatorname{dist}_B(v,p_w)-w\ge0\) and
\(b_B=\max_v\operatorname{dist}_B(v,p_0)\ge w\), including path vertices
in both maxima. In addition to \(a_{B}+w\le 2s\) and \(b_{B}\le 2s\),
one has \(a_{B}+b_{B}\le 2s\). If \(a_{B}=0\) or the maximum for
\(b_{B}\) is attained on the path this follows from the preceding
bounds. Otherwise let C,D contain maximizing vertices for
\(a_{B},b_{B}\). Then \(a_{B}\le s_{C}-\beta _{C}\) and
\(b_{B}\le s_{D}+\alpha _{D}\). For \(C=D\) use
\(\alpha _{C}\le \beta _{C}\); for \(C\ne D\) use
\(\alpha _{D}-\beta _{C}\le 0\) or the covered-gap bound
\(2(s-s_{C}-s_{D})\).

Keep each block internally fixed, in its original order, and let its
starting coordinate \(t_{i}\) vary. Unlike the original decomposition,
allow consecutive blocks to touch at one P vertex. The constraints are
\(a_{i}\le t_{i}\le d-b_{i}\), and t\_\{i+1\}\textgreater=t\_i+w\_i. The
internal diameters are unchanged. Joining at a shared endpoint creates
no new shortcuts on P, and distances between distinct blocks factor
through their intervening P segment. For x in a block to the left of y's
block,
\(\operatorname{dist}(x,y)=\operatorname{dist}(x,p_{d})+\operatorname{dist}(y,p_{0})-d\);
hence the two endpoint bounds suffice for all cross-block diameter
constraints. Path vertices between blocks also obey these bounds. All
translated graphs have the same n,m,d.

For each off-path vertex x in block i the sum of its distances to P is a
quadratic in \(t_{i}\) with leading coefficient 1. The left and right
tails contain \(t_{i}\) and \(d-t_{i}-w_{i}\) path vertices
respectively. Each tail contributes an arithmetic progression with
quadratic term one half the square of its length; the internal
contribution is fixed. Distances between different blocks are affine in
their positions, and internal distances are constant. Thus translating
any nonempty chain of blocks together gives a strictly convex quadratic.
At a maximizer every maximal touching chain must be pinned to an
individual lower or upper bound: otherwise both integer shifts by one
remain feasible and one strictly increases W.

A chain pinned at a lower bound has its rightmost coordinate at most 2
times its total off-path vertex count, by \(a_{i}+w_{i}\le 2s_{i}\) and
\(w_{j}\le 2s_{j}\) for subsequent blocks. The reflected statement holds
for upper-pinned chains. For \(d>2k\) no chain can be pinned both ways.
Indeed, if a chain has a lower pin i and an upper pin j with \(i\le j\),
then \(d=a_i+b_j+\sum_{u=i}^{j-1}w_u\). For \(i<j\) this is at most
twice the chain size by \(a_{i}+w_{i}\le 2s_{i}\) and
\(b_{j}\le 2s_{j}\); for \(i=j\) use \(a_{i}+b_{i}\le 2s_{i}\). For
\(j<i\) the intervening widths are subtracted and
\(a_{i}+b_{j}\le 2(s_{i}+s_{j})\). Thus the chains split into separated
left/right groups. This argument requires the strict inequality
\(d>2k\): at \(d=2k\) a path of k edges attached at the midpoint of P
can be doubly pinned. Let L,R be their total off-path counts. Each
group's attachment intervals lie within 2L or 2R coordinates of its
endpoint, respectively. Because \(2L+2R=2k<d\), the two sides are
separated, have no edges between their off-path vertices, and cross-side
distances are \(d-e_{x}-e_{y}\) (with e defined in Section 3). Every
pair of off-path vertices in a group has distance \(\le 2\) times that
group's size. Within a block this was proved above. In different blocks
of a touching chain it follows by routing through their path boundaries
and the intervening widths, all bounded by twice the contributing
counts. In different left-pinned chains, route via p0: every vertex of a
chain of size s' has distance to p0\textless=2s'. For a vertex in the
pinned block use \(a_{i}+b_{i}\le 2s_{i}\); before the pin use
\(t_{j}\le a_{i}\le 2s_{i}\) and \(b_{j}\le 2s_{j}\); after the pin use
\(a_{i}+w_{i}\le 2s_{i}\) plus intervening widths and
\(b_{j}\le 2s_{j}\). Reflect for right chains. This proves the asserted
group diameter bound.

\section{3. Exact identity, capacity deficit and deficit
one}\label{exact-identity-capacity-deficit-and-deficit-one}

\subsection{3.1 Identity and zero
deficit}\label{identity-and-zero-deficit}

Orient coordinates locally from each endpoint. For an off-path vertex
\(x\), put
\(e_x=d-\operatorname{dist}(x,\text{opposite endpoint})\ge0\), let
\(a_x\) count its direct neighbors on \(P\), and put
\(\sigma_x=e_x+2-a_x\). For any such neighbor \(p_j\), the triangle
inequality gives
\(d-j-1\le\operatorname{dist}(x,\text{opposite endpoint})\le d-j+1\).
Hence \(j\in\{e_x-1,e_x,e_x+1\}\), with only \(0,1\) possible at
\(e_x=0\). It follows that \(a_x\le e_x+2\) and \(\sigma_x\ge0\).

Let \(H\) count missing same-side off-path edges, \(S=\sum_x\sigma_x\),
\(E_L=\sum_{x\in L}e_x\), \(E_R=\sum_{x\in R}e_x\) (using \(L,R\) also
for the vertex groups), \(E=E_L+E_R\), and \(\delta=H+S\). The number of
path attachments is \(E+2k-S\), and the off-path edge count is
\(\binom{k}{2}-LR-H\), because the sides have no cross edges. Thus

\[r=E+A-LR-\delta,\qquad \delta=h-LR+E.\]

Define the profile and loss terms by

\[T_x=\sum_{j=0}^d\bigl(\operatorname{dist}(x,p_j)-j+e_x\bigr),\qquad U_x=T_x-2e_x-1,\]
\[U=\sum_x U_x,\qquad J=\sum_{\{x,y\}\text{ same side}}(\operatorname{dist}(x,y)-1).\]

Here the path coordinates in the profile are local to the side. Summing
the path-path, path-off and off-off distances, with cross distances
\(d-e_x-e_y\), gives exactly

\[W=F+hd+r-\delta d+\delta+U+J-R E_L-L E_R.\]

For clarity, the path-off distance sum is
\(T_x+\binom{d+1}{2}-(d+1)e_x\). Same-side pairs contribute
\(\binom L2+\binom R2+J\); cross-side pairs contribute
\(LRd-RE_L-LE_R\). Substitution of \(\sum_xT_x=U+2E+k\) and
\(E=\delta-h+LR\) gives the displayed identity.

If \(\sigma _{x}=0\), \(e=0\) or 1 and the P-neighbor mask is exactly
\{0,1\} or \{0,1,2\}, respectively. For \(e=0\), distances to p0 and
\(p_{j}\) \((j\ge 1)\) are 1 and j. For \(e=1\), distances to p0,p1 are
1 and 1, and those to \(p_{j}\) for \(j\ge 2\) are j-1. The lower bound
\(\operatorname{dist}(x,p_{j})\ge j-e_{x}\) follows from the
opposite-endpoint triangle inequality. These profiles give
\(T_{x}=2e_{x}+1\).

If \(\delta =0\), both sides are cliques and all masks are good. By
reflection take \(L\ge R\). \(E=LR-h\ge 0\) forces \(R\ge R_0\). The
constant loss is R \(E_{L}+L\) \(E_{R}\ge R(LR-h)\ge R_0\) q, with
strict inequality if \(R>R_0\) since j(k-j) increases on the allowed
integer interval. At \(R=R_0\) equality requires all q shortened
vertices on the larger side; if \(q=0\) there are none. In the strict
middle range \(L_0>R_0\): for even \(k=2b\), \(h\le b^2-1=(b-1)(b+1)\),
and for odd k the groups have unequal sizes. By minimality of R0,
\((R_0-1)(L_0+1)<h\), so \(0\le q\le L_0-R_0<L_0\). In particular
\(qR0\le R_0(k-2R0)\le k^2/8\), and by integrality \(qR0\le M\). This is
exactly \(H_*\), up to reflection and vertex permutations.

\subsection{3.2 Deficit one}\label{deficit-one}

If \(\delta =1\) and \(H=1\), all masks are good, \(U=0\) and \(J=1\)
(both endpoints of the missing edge share the same p1). If \(S=1,H=0\),
there is just one bad vertex and the side is a clique, so \(J=0\). Its
possible (e,a) are (0,1),(1,2),(2,3). For (0,1), a direct neighbor p1
gives \(U_{x}\le 1\); if the sole neighbor is p0, the opposite-endpoint
geodesic must next use a good off-path vertex with \(e=1\), which gives
the same bound. For (1,2), p2 must be a neighbor, since no other
off-path vertex has \(e=2\); the two masks \{0,2\},\{1,2\} each give
\(U_{x}\le 1\). For (2,3) the mask is \{1,2,3\}, giving \(U_{x}=2\); the
opposite group is nonempty since otherwise \(\delta =h+E\ge 3\), using
\(h>0\) from the strict middle range. Its cross loss is \(\ge 2\). In
all cases \(\delta +U+J-R\) \(E_{L}-L\) \(E_{R}\le 2\). Thus
\(\delta =1\) loses strictly to \(H_*\) when \(d\ge M+3\) because
\(qR0\le M\).

\section{4. Analytic proof for k \textgreater=
19}\label{analytic-proof-for-k-19}

\subsection{4.1 Oriented charges}\label{oriented-charges}

Use the notation of Section 3. Within a group of s off-path vertices,
for each x take a shortest path to the opposite endpoint, up to its
first P vertex at local coordinate \(t_{x}\). Let \(\ell _{x}\) be its
length to P, \(c_{x}=\operatorname{max}(\ell _{x}-2,0)\), and
\(z_{x}=\sigma _{x}+c_{x}\). Every off-path step of this path increases
\(e\) by exactly one. The prefix followed by the path from \(p_{t_x}\)
to the opposite endpoint is shortest, so \(d-e_x=\ell_x+d-t_x\), or
\(t_x=e_x+\ell_x\). The vertices at steps \(2,...,\ell _{x}-1\) are
nonadjacent to x, because the path is shortest. Charge their missing
unordered pairs to x. Each pair can be charged at most once: the charge
goes from its smaller e to its larger e. Let \(C=\sum_x c_x\). Then
\(C\le H\). Both ends of a charged pair have positive \(\sigma\): the
source has \(\ell \ge 3\), whereas \(\sigma =0\) supplies a direct
optimal P neighbor and no shortest path can be assumed to have
\(\ell \ge 3\) if we choose shortest \(\ell\) among shortest paths. We
henceforth require a minimum-length prefix to P among shortest paths to
the opposite endpoint. In particular, no earlier vertex of that prefix
can have an optimal P neighbor: using it would give a shorter prefix. An
optimal P neighbor is at coordinate e+1. This rule is essential;
choosing a longest prefix is not valid. The target has \(e\ge 2\) and
\(\sigma =e+2-a\ge e-1\ge 1\). Thus \(C\le \operatorname{binom}(S,2)\),
where \(S=\sum_x\sigma_x\) (the number of positive \(\sigma\) vertices
is at most S). The chosen paths remain inside their group until reaching
P. In particular \(\ell \le s\) and \(t\le 2s\).

\subsection{4.2 Individual profile
bounds}\label{individual-profile-bounds}

Write \(e=e_{x}\), \(a=a_{x}\), \(\ell =\ell _{x}\), \(c=c_{x}\),
\(z=z_{x}\), \(t=t_{x}\) and \(U_{x}=T_{x}-2e-1\). We prove
\(U_{x}\le 2sz\) and \(U_{x}\le z(z+1)\). If \(\sigma =0\),
\(U_{x}=c=0\) by the good mask and the shortest-ell choice. In all other
cases \(t\le z+2\).

For \(a=0\), \(\sigma =e+2\), \(\ell \ge 2\), and \(t=e+\ell =z\). For
\(j\ge t\), the route via \(p_t\) and the opposite-endpoint lower bound
give \(\operatorname{dist}(x,p_j)=j-e\). For \(j<t\), the same route
bounds the profile term by \(2(t-j)\). Consequently
\(T_x\le\sum_{j=0}^{t-1}2(t-j)=t(t+1)\). If \(t\le 2s-1\),
\(U_{x}\le 2st\); if \(t=2s\), use \(e=t-\ell\) and \(\ell \le s\) to
obtain \(U_{x}\le t(t-1)+2\ell -1<t^2=2st\). Also \(U_{x}\le z(z+1)\).

For \(a\ge 1\) let b be the largest P-neighbor coordinate and
\(g=e+1-b\). All P neighbors lie in \{e-1,e,e+1\} with nonnegative
coordinates, so \(0\le g\le 3-a\). For \(j\le b\) the route via
\(p_{b}\) bounds the profile term by e+1+b-2j; for \(b<j<t\) it bounds
it by g; for \(j\ge t\) it vanishes. Summing gives

\[U_x\le e^2+e+1+g(\ell+g-e-3)
\le e^2+e+1+g(c+g-e-1).\]

\begin{itemize}
\tightlist
\item
  \(a=1\): \(\sigma =e+1\); \(g\le 2\) (g\textless=1 if e=0), so
  \(g(g-e-1)\le 0\). Hence \(U_{x}\le \sigma ^2+2c\).
\item
  \(a=2\): \(\sigma =e\), and the bad case has \(e\ge 1\). If \(b=e\),
  the bound gives \(U_{x}\le e^2+1+c\le e(e+1)+c\). If \(b=e+1\), the
  second neighbor is e or e-1; if the mask is \{e,e+1\}, the terms for
  \(j\le e\) are at most 2(e-j)+1, and the later terms vanish, so
  \(T_{x}\le (e+1)^2\). If the mask is \{e-1,e+1\}, the terms for
  \(j\le e-1\) are at most 2(e-j), the term at \(j=e\) is at most 2, and
  later terms vanish. Hence \(T_{x}\le e(e+1)+2\). Both bounds imply
  \(U_{x}\le e^2+1\le e(e+1)\) for \(e\ge 1\).
\item
  \(a=3\): the full three-neighbor mask yields \(T_{x}\le e(e+1)+1\),
  hence \(U_{x}\le e(e-1)=\sigma (\sigma +1)\).
\end{itemize}

In these cases \(\sigma \le 2s\) and, for \(a=2,3\),
\(\sigma +1\le 2s\). Each stated bound is at most 2s z and z(z+1), since
\(\sigma \ge 1\). The relation \(t=e+\ell \le z+2\) holds for \(a=1,2\);
for \(a=3\) the direct optimal neighbor gives \(\ell =1\),
\(t=\sigma +2\). For \(\sigma =0\) it holds as well.

\subsection{4.3 Group bounds}\label{group-bounds}

For a charged pair at step i of a chosen path, its distance is exactly
i, so its contribution to J is i-1. Total charged \(J\) is
\(\sum_x c_x(c_x+1)/2\le (s-1)C/2\) because \(c_{x}\le s-2\) whenever
\(c_{x}>0\). Each uncharged missing pair contributes at most 2s-1 by
localization. Since \(\sum_x z_x=S+C\), within a group

\[\delta+U+J\le(2s+1)S+2sH+\frac{(s+1)C}{2}\le\frac{(5s+1)\delta}{2}.\tag{B1'}\]

An independent estimate uses \(z\le \delta\) and
\(U\le\sum_x z_x(z_x+1)\le(\delta+1)(S+C)\). For any missing pair, the
first P coordinates give distance
\(\le 2max(t_{x},t_{y})\le 2\delta +4\). Charged J is at most \(\delta\)
C (indeed \((\delta+1)C/2\) for C\textgreater0). Therefore

\[\delta+U+J\le(\delta+2)S+(2\delta+4)H-2C\le(2\delta+4)\delta.\tag{B2'}\]

Both inequalities hold after summing over groups with total k and total
\(\delta\), by replacing each s and each group \(\delta\) by their
global bounds. Explicitly, the quadratic bound sums safely because
\(\sum_g(2\delta_g+4)\delta_g\le(2\delta+4)\delta\); the loss function
\(\phi (t)=t(d-2t-4)\) satisfies
\(\phi (a)+\phi (b)-\phi (a+b)=4ab\ge 0\). The sharper expanded B1' also
holds globally:

\[\delta+U+J\le(2k+1)S+2kH+\frac{(k+1)C}{2}.\tag{B3'}\]

\subsection{4.4 Exclusion of positive
deficits}\label{exclusion-of-positive-deficits}

Let \(M=\operatorname{floor}(k^2/8)\), \(d\ge M+3\), and recall
\(qR0\le M\) and the exact identity in Section 3. For \(k\ge 20\) one
has \(d>2k\). Deficit 0 is classified in Section 3 and deficit 1 is
excluded by the deficit-one argument in Section 3. For
\(2\le \delta \le 20\), B2' gives a loss \(\ge \delta (d-2\delta -4)\);
it exceeds M at both endpoints when \(k=20\) (90 and 180 versus 50).
Concavity in \(\delta\) and monotonicity of the margins in M prove it on
the full interval. For \(\delta \ge 21\), B1' gives loss
\(\ge 21(d-(5k+1)/2)\), which is \(52.5>50\) at \(k=20\). Its margin
over M grows since M increases by at least 5 per unit k and
\(20*5>21*5/2\). Thus all positive deficits lose strictly. For \(k=19\),
\(M=45,d\ge 48>2k\). B2' excludes \(2\le \delta \le 20\) (both endpoint
losses are 80). For \(\delta \ge 21\) write \(B=H-C\ge 0\). By B3',
\(d*\delta -(\delta +U+J)\ge 9S+10B\). If this were \(\le 45\),
\(S+B\le 5\), so \(C\le \operatorname{binom}(S,2)\le 10\) and
\(\delta =S+B+C\le 15\), a contradiction. Thus this loss too is strictly
greater than M. The nonnegative cross losses can only help. This proves
the upper threshold for \(k\ge 19\). The matching lower threshold is
proved in Section 7.

\section{5. Integer certificate for 11 \textless= k \textless=
18}\label{integer-certificate-for-11-k-18}

\subsection{5.1 Density-controlled
localization}\label{density-controlled-localization}

For this section write \(d_0=M+3\) and
\(r_{\min}=A-\lfloor k^2/4\rfloor+1\). All admissible middle-range
graphs have \(r\ge r_{\min}\). We use \(d\ge d_0\). In Section 5,
variables \(H,S,C,\delta,U,J\) in a single-side calculation denote their
restrictions to that side; sums over both sides restore the global
quantities.

Let the off-P components have sizes \(u_{i}\) and intrinsic diameters
\(D_{i}\) in the off-P graph, and let c be their count. Write
\(B^*=\sum_i(D_i+2)\le k+c\). A component's attachment interval has
width \(\le D_{i}+2\), and each of its vertices can reach either extreme
attachment in \(\le D_{i}+1\) steps. In each merged block use
\(B_{\mathrm{block}}=\sum_i(D_i+2)\) over its components. The width w,
a+w and b (notation of Section 2) are \(\le B_{block}\), by routing via
an extreme attachment and bounding the residual gap by the other
component interval widths. For a+b, if the two maximizing vertices are
in distinct off-P components, use \(a\le D_{i}+1-\beta _{i}\) and
\(b\le D_{j}+1+\alpha _{j}\) in block coordinates; a positive gap
\(\alpha _{j}-\beta _{i}\) is covered by the remaining intervals. Hence
\(a+b\le B_{block}\). Path-vertex maxima obey the same bound. If the
vertices are in the same component,
\(a+b\le 2(D_{i}+1)+\alpha _{i}-\beta _{i}\le 2(D_{i}+1)\).

The convex touching-chain argument of Section 2 is unchanged. Two pins
in different blocks would imply \(d\le B^*\) (sum the intervening
widths; if the upper pin precedes the lower pin, subtract intervening
widths). Two pins in the same block require d\textless=max(B*,2
max\_i(D\_i+1)). If these possibilities are excluded, all chains are
one-side pinned, and each side has attachment coordinates \(\le\) its
own sum \(B=\sum_i(D_i+2)\), by the lower-pin bound \(a+w\le B_{block}\)
and subsequent block widths. \(B^*<d\) separates the two sides. The
Section 2 within-group bound 2s continues to hold.

Every off-P vertex has at most three P neighbors, so the number of off-P
edges is \textgreater=rmin-2k. For c components, it is
\(\le \operatorname{binom}(k-c+1,2)\). For a component of size u and
diameter D, a diametral path in that component has D edges; each of its
u-D-1 remaining vertices has at most three neighbors on the path, hence
the component has at most
\(f(u,D)=D+3(u-D-1)+\operatorname{binom}(u-D-1,2)\) edges. Thus
f(u,D)+binom(k-u,2)\textgreater=rmin-2k for any component. The integer
certificate computes upper bounds on c and D using these necessary
conditions and checks \(d_0>k+c\) and \(d_0>2(D+1)\), except \(k=11\)
where \(d_0=2(D+1)=18\). In this equality case double pins would require
a single component of diameter 8 whose attachment interval has width
zero (alpha=beta, by the same-component inequality). Its P-neighbor
count is then \(\le u\), so
\(r\le f(u,8)+\operatorname{binom}(k-u,2)+2(k-u)\). For every
\(9\le u\le 11\) this is \textless rmin. Hence double pins are
impossible there too. This establishes localization for all
\(11\le k\le 18\) and \(d\ge d_0\).

\subsection{5.2 Relaxation}\label{relaxation}

For a side of size s, enumerate all integer partitions \((u_{i})\) of s.
Let \(H_c=\binom{s}{2}-\sum_i\binom{u_i}{2}\) be the number of
cross-component missing off-P pairs; \(H_{i}\) the count within
components; \(H=H_{c}+H_{i}\). Connectivity yields
\(0\le H_i\le\sum_i\binom{u_i-1}{2}\). Also
H\textless=2k+binom(s,2)+binom(k-s,2)-rmin by edge counting, allowing
the other side to have all its off-P edges. For each component, a
possible intrinsic diameter v must satisfy both
\(\operatorname{binom}(v,2)\le H_{i}\) (nonconsecutive pairs of a
diametral geodesic) and \(\operatorname{binom}(u,2)-f(u,v)\le H_{i}\).
Let D\_i' be the maximum such integer \(0\le v<u\),
\(D=\max_iD_i\prime\) and \(B=\sum_i(D_i\prime+2)\). These are upper
bounds, not claimed realizable jointly. Attachment coordinates
\(t_{x}\le B\) and geodesic prefix lengths \(\ell _{x}\le D+1\). Since
\(e_{x}=t_{x}-\ell _{x}\le B-1\) and each component has a P neighbor,
\(0\le S=\sum_x\sigma_x\le(B+1)s-c\). Put \(\delta =S+H\).

\subsection{5.3 Profile and pair bounds}\label{profile-and-pair-bounds}

As in Section 4 let \(c_{x}=\operatorname{max}(\ell _{x}-2,0)\),
\(z_{x}=\sigma _{x}+c_{x}\), \(C=\sum_x c_x\). Choose minimum prefix
length among shortest paths to the opposite endpoint. For \(a_{x}\ge 1\)
the individual bounds of Section 4 yield \(U_{x}\le B\) \(z_{x}\): for
\(a=1\) use \(\sigma \le B\) and \(U_{x}\le \sigma ^2+2c_{x}\); for
\(a=2\) use \(\sigma +1\le B\) and
\(U_{x}\le \sigma (\sigma +1)+c_{x}\); for \(a=3\) use
\(\sigma +1\le B\). For \(a=0\), \(t=z\) and \(U_{x}\le t(t+1)-2e-1\).
If \(t<B\) this is \(\le Bz\); if \(t=B\) and \(2D+1\le B\), substitute
\(e=t-\ell\) to obtain \(U_{x}\le t^2-t+2\ell -1\le t^2-t+2D+1\le t^2\),
since \(t=B\). Otherwise \(U_{x}\le (B+1)z\). Thus \(U\le K(S+C)\) where
\(K=B+\mathbf{1}_{\{2D+1>B\}}\). Also
\(U\le\sum_x z_x(z_x+1)\le(S+C)(S+C+1)\).

For a missing pair in different off-P components of the same side, route
through the chosen P coordinates. Its distance is
\(\le \ell _{x}+\ell _{y}+|t_{x}-t_{y}|\le B+D+1\), as the earlier
coordinate minus its own prefix length is nonnegative. Combine with 2s.
Within the same component distance\textless=D. A charged pair at step i
contributes i-1 to J, so charged \(J\le DC/2\), since each
\(c_{x}\le D-1\). For \(D\ge 2\), therefore

\[J\le(\min(2s,B+D+1)-1)H_c+(D-1)(H_i-C)+\frac{DC}{2}.\]

For \(D<2\) there are no actual within-component missing pairs or
charges. For relaxation use max(1,D) as coefficient of \(H_{i}\) in
\(\delta +J\). These estimates give the doubled loss \texttt{loss1} in
the script, i.e.~a lower bound for 2{[}d delta-(delta+U+J){]}.
Independently, bound B2' of Section 4 yields loss2. And
\(t_{x}\le z_{x}+2\le S+C+2\) gives \(J\le (2(S+C)+3)H\), yielding
loss3. The maximum of these lower bounds is valid. For side
\(\delta =1\) the Section 3 deficit-one cases give \(\delta +U+J\le 3\)
even without the cross loss; hence the additional lower bound 2(d-3).
Side \(\delta =0\) has zero such loss.

\subsection{5.4 Sharper charge bound}\label{sharper-charge-bound}

All charged pairs lie inside an off-P component; \(C\le H_{i}\),
\(C\le \operatorname{binom}(S,2)\), \(C\le\sum_i u_i\max(u_i-2,0)\). Let
\(j=\max_x c_x\ge1\) and p be the number of vertices of positive
\(\sigma\). A source with \(c=j\) has j+2 off-P vertices on its geodesic
prefix (at steps 0,\ldots,j+1). Each has positive \(\sigma\): before the
last, an optimal direct P neighbor would give a shorter prefix, so the
only possible P neighbors are at e-1 and e. For \(e\ge 1\), \(a\le 2\)
gives \(\sigma \ge e\ge 1\); for \(e=0\) only coordinate 0 is possible,
giving \(a\le 1\) and \(\sigma \ge 1\); at the last \(e\ge j+1\ge 2\),
hence \(\sigma \ge j\). Along the prefix e increases by one each step.
At steps 2,\ldots,j the excess of \(\sigma\) over 1 is at least i-1; at
the last it is at least j-1. Thus \(p\le (S-(j-1)(j+2)/2)\),
\(p\ge j+2\). Also \(p\le s\). In a list of the p positive-sigma
vertices sorted by e, a source at rank i (1-based) with \(c_{x}>0\) has
at least \(c_{x}+1\) later vertices, so
\(c_{x}\le \operatorname{min}(j,\operatorname{max}(p-i-1,0))\). Sum over
ranks and maximize over \(1\le j\le \operatorname{min}(s-2,D-1)\),
allowing zero. This is \texttt{charge\_cap}. The script substitutes an
upper bound for C in all loss bounds, which are nonincreasing in C. No
assumption that the relaxed integer data can be realized by a graph is
needed.

\subsection{5.5 Certificate}\label{certificate}

For reproducibility, the three doubled loss bounds used in the program
are

\[\begin{aligned}
\mathrm{loss}_1={}&2[(d-K-1)S+(d-\min(2s,B+D+1))H_c\\
&+(d-\max(1,D))H_i]-(2K+2-D)C,\\
\mathrm{loss}_2={}&2[d\delta-(2\delta+4)\delta],\\
\mathrm{loss}_3={}&2[d\delta-\delta-(S+C)(S+C+1)\\
&-(2(S+C)+3)H].
\end{aligned}\]

In loss1 omit the last term when \(D<2\) (actual C=0). The upper bound
on C is safe since loss1 and loss3 are nonincreasing in C; loss2 is
independent of C. For an empty group the table is \{0:0\}. For
\(\delta =1\) also use 2(d-3), and for \(\delta =0\) use zero.

The script records the minimum doubled lower loss for each side size and
\(\delta\) over all the above necessary integer data. Pair the two sides
\((L\ge R)\), omitting total \(\delta <2\). For \(k=11..18\), every
paired doubled lower loss is strictly greater than 2M. The cross loss R
\(E_{L}+L\) \(E_{R}\) is nonnegative and may be ignored for these cases.
Thus the exact identity of Section 3 excludes all \(\delta \ge 2\) at d0
and all larger d, since every doubled loss has coefficient 2
\(\delta _{g}\) in d (zero for the zero-deficit state and positive
otherwise). Total \(\delta =1\) is excluded by Section 3 (including
cross loss), total \(\delta =0\) classified by Section 3. The matching
lower-threshold rival is Section 7. This closes \(k=11..18\) using the
analytic lemmas and the finite arithmetic certificate. The sensitivity
analysis in the independent stage-4 audit found that \(k=11\) depends
essentially on charge\_cap; for \(k=12..18\) loss1 alone suffices. The
\(k=11\) localization equality case is handled in Section 5.1. The
script also reports \(k=10\) as an explicitly unsuccessful diagnostic of
the relaxation (localization not certified for that case). Its remaining
integer states are not graph counterexamples.

\section{6. Exact finite certificate for 4 \textless= k \textless=
10}\label{exact-finite-certificate-for-4-k-10}

\subsection{6.1 Finite-diameter
reduction}\label{finite-diameter-reduction}

Let \(M=\operatorname{floor}(k^2/8)\), \(d_0=M+3\),
\(A=\operatorname{binom}(k+1,2)\), \(h=A-r\),
\(0<h<\operatorname{floor}(k^2/4)\). Let \(H_*=H_{d,r}\) be the
canonical candidate of Section 1, with \(W(H_*)=F(d,k)+hd+r-qR_0\).
Suppose \(d_0\le 2k\) (for \(k\ge 4\) this is exactly k\textless=14). We
use this reduction for \(k=4..10\) only. To certify the exact uniform
threshold d0 it suffices to certify the value and equality structure for
every integer d in {[}d0,2k{]}, and exhibit a noncanonical rival at d0-1
as in Section 7.

Proof. Consider a maximizer G of diameter \(d>2k\) and fix a diametral
geodesic P. By Section 2, its left and right groups contain L,R off-path
vertices with \(L+R=k\). All attachment coordinates from the respective
ends are at most 2L,2R; all distances within each group (including the
off-path vertices' distances to its near endpoint) are at most 2L,2R.
The local fixed blocks are unchanged if the empty central segment is
shortened so that the path has length \(D=2k\). Specifically keep the
first 2L and last 2R edges of P and identify their inner endpoints after
deleting the intervening path. If \(L=0\) or \(R=0\), the corresponding
kept segment consists of the endpoint alone. The two segments may meet
at exactly one vertex; they share no edge, and there are no edges
between their off-path sets. This operation removes exactly d-D vertices
and d-D edges; k and r are unchanged. No loops or duplicate edges are
created. Each side is attached to the other at the shared path vertex,
so within-side distances are preserved. The retained path remains
geodesic, since no shortcut existed in either original side. For an
off-path vertex x, its distance to the opposite endpoint is now
\(D-e_{x}\), with \(e_{x}\) as before. This is at most D. Distances to
path vertices within its own segment are at most D (indeed \(\le 2L\) or
2R from the block estimates, or by the fact that the maximum of the
distance profile on that segment is bounded by the corresponding block
diameter and endpoint bounds); more directly the original within-side
paths have length at most 2s for any two vertices of the kept side: for
a vertex x in a block, distances to path vertices outside that block are
maximized at one of the segment endpoints, both at distance \(\le 2s\),
and distances inside the block are \(\le 2s\) by Section 2. To justify
the distance to the inner endpoint at coordinate 2s, use the
opposite-endpoint route formula, which gives \(2s-e_{x}\). Distances to
path vertices on the other side are \(\le D-e_{x}\). Cross-group
off-path distances are \(D-e_{x}-e_{y}\), also \(\le D\). Thus the new
graph G' has diameter D exactly.

The capacity deficit \(\delta =H+sum\) \(\sigma _{x}\) is nonnegative
and unchanged. The U and J profile terms and the cross-loss R
\(E_{L}+L\) \(E_{R}\) of the exact identity in Section 3 are unchanged:
the deleted path coordinates lie after every left attachment and before
every right attachment, where the oriented profile is zero. This gives

\[[W(G)-W(H_{d,r})]-[W(G\prime)-W(H_{D,r})]=-\delta(d-D).\]

If the diameter-D certificate says G' cannot exceed \(H_{D,r}\), neither
can G. Equality at \(d>D\) requires \(\delta =0\), and the deficit-zero
classification of Section 3 gives exactly \(H_{d,r}\) up to isomorphism.
This proves the reduction for all larger d.~The finite interval is not
an extrapolation from a sampled sequence of d: it follows from the exact
identity and the localization lemma. The reduction uses the strict
inequality \(d>2k\); the anchor \(d=2k\) itself is verified by
enumeration, not by assuming localization there.

\subsection{6.2 Enumeration universe}\label{enumeration-universe}

Fix an ordered diametral geodesic P of length d.~Any off-path vertex has
a P-neighbor mask of span at most two, or the empty mask. Relabel
off-path vertices to order the masks nondecreasingly. For each resulting
tuple enumerate all subsets of the possible off-path edges. Reject an
edge between vertices with nonempty masks if some pair of their P
neighbors differs in coordinate by more than 3: the edge would give a
shortcut of length 3. Empty masks are kept, so vertices not directly
attached to P are included. The implementation uses uint64\_t for both
the adjacency masks and off-path edge subsets, sufficient for
\(k\le 10,d\le 20\) (45 off-path edge bits, \(\le 31\) vertices). The
non-isomorphic graphs are not counted; duplicate labeled representations
may occur.

\subsection{6.3 Basic and prefix bounds}\label{basic-and-prefix-bounds}

For a partial graph K consisting of P, all fixed P-neighbor edges, and a
subset of off-path edges, let
\(b(u,v)=\operatorname{min}(d,dist_{K}(u,v))\), with infinity capped at
d.~Every admissible completion G contains K and has diameter d, so
\(dist_{G}(u,v)\le b(u,v)\). If an undecided edge uv is added, its own
pair has distance exactly one; this reduces its contribution to the
upper bound by b(u,v)-1. Different added edges refer to distinct
unordered pairs. Hence for t more edges, \(W(G)\le sum\) b - sum of the
t smallest costs b(u,v)-1 among the remaining edges. The actual
reduction on other pairs is nonnegative, so ignoring it is safe. Apply
this bound separately to every feasible final \(r=|E(K)|-(n-1)+t\) in
the strict middle interval. Prune only when every such upper bound is
strictly smaller than the candidate value, or when a path shortcut or
impossible edge count excludes every completion. Equality is never
pruned by an objective bound. At the initial mask stage the capped
distances can be computed directly using only P; the off-path pair upper
bound is \(\operatorname{min}(d,min_{j}[b(x,p_{j})+b(y,p_{j})])\). If
one or both masks are empty this gives d.~This is an upper bound even in
disconnected partial graphs. No approximation or statistical selection
is used.

The prefix bound further restricts the search while choosing masks. For
a completed graph let c be the total number of P-neighbor edges, and e
the number of off-path edges, so \(r=c+e-k\). For each mask m write
a(m)=popcount(m), and \(u(m)=sum_{j}\) min(d,1+min\_\{i in
m\}\textbar i-j\textbar), setting u(empty)=(d+1)d.~Then

\[W(G)+(d-1)r\le W(P)+\sum_m[u(m)+(d-1)a(m)]+d\binom{k}{2}-(d-1)k.\]

Indeed each off-path pair contributes at most d after adding d-1 to its
distance if it is an edge (then its distance is one). A partial mask
tuple gives an upper bound on the RHS by replacing every undecided unary
term with the maximum over all mask indices at or above the current
index. If this is strictly below min\_r{[}W(H\_\{d,r\})+(d-1)r{]},
discard the entire prefix. This is again a relaxation and equality is
retained.

\subsection{6.4 Symmetry and stronger prefix
bounds}\label{symmetry-and-stronger-prefix-bounds}

Fix the sorted tuple of P-neighbor masks as in Section 6. Process
off-path vertices in label order. At row i, all edges incident to an
earlier off-path vertex have been decided; edges between unprocessed
vertices are undecided except for mask-incompatible edges that can never
occur. Partition the possible neighbors \(j>i\) by their P-neighbor mask
and their vector of adjacencies to vertices before i. A permutation
within a class preserves the masks, all earlier decisions, and
compatibility of every undecided edge. It fixes i. Therefore each subset
of s equivalent candidate neighbors is equivalent, under a relabeling of
the not-yet-processed vertices, to a prefix of that class of the same
cardinality. Test the s+1 prefix lengths, rather than \(2^s\) subsets.
This does not assert that the resulting complete graphs are all
nonisomorphic; it does guarantee at least one representation of every
admissible completion up to relabeling. Canonical shape recognition is
invariant under such off-path relabeling. Each row's edges become
decided before advancing to the next row. The distance/edge-count
pruning of Section 6 remains unchanged and retains equality.

Let \(m_{i}\) be the P-neighbor mask, a(m) its population count, u(m)
the capped distance sum to P via the mask, and p(m,n) the capped pair
distance upper bound computed via P, as in Section 6. An incompatible
pair (as defined in Section 6 by the length-three shortcut test) cannot
be an edge and hence contributes at most p(m,n) to W+(d-1)r. Every other
pair contributes at most d, whether it is an edge or not. Put
penalty(m,n)=d-p(m,n) for an incompatible pair and zero otherwise. Then

\[\begin{aligned}W+(d-1)r\le{}& W(P)+\sum_i[u(m_i)+(d-1)a(m_i)]\\&+d\binom{k}{2}-(d-1)k-\sum_{i<j}\operatorname{penalty}(m_i,m_j).\end{aligned}\]

For a mask prefix, retain all fixed unary terms and fixed-fixed
penalties. For each undecided mask bound its unary term minus all its
fixed-prefix penalties by the maximum over allowable suffix masks, and
discard all nonnegative undecided-undecided penalties. This is an upper
bound. Also, the maximum achievable edge count is bounded by the
compatible pairs among fixed masks plus all pairs involving undecided
masks and at most three P attachments per undecided vertex. The code
discards prefixes that cannot meet the strict middle edge range or
cannot reach the minimum, over that range, of candidate W+(d-1)r. Only
strict objective deficits are pruned.

\subsection{6.5 Connectivity}\label{connectivity}

Let z be the number of vertices with empty P-neighbor mask. Every
admissible connected graph has at least one nonempty mask. For any
empty-mask vertex x, a shortest path in the off-path graph to the set of
nonempty-mask vertices has length at most z: all vertices before the
final nonempty-mask vertex have empty masks and are distinct. If the
final vertex has mask m, then
\(\operatorname{dist}(x,p_{j})\le \operatorname{min}(d,z+md(m,j))\),
where md is the capped one-edge-plus-path distance of Section 6. Define

\[U_h=\max_{m\ne\varnothing}\sum_{j=0}^d\min(d,h+\operatorname{md}(m,j)).\]

The total empty-mask contribution to the P-off distance sum is at most z
\(U_{z}\), replacing the previous z(d+1)d.~The sorted mask order places
the empty mask first. Once a nonempty mask is chosen, z is fixed for
every completion of the prefix; deduct z(d+1)d-z \(U_{z}\) from the
prefix bound. The preliminary connectivity bound uses this deduction,
retaining the earlier pair penalties (all empty-mask pair penalties are
zero).

A further safe strengthening uses the distances \(h_{i}\) of the
empty-mask vertices to the set of nonempty-mask vertices, arranged in
increasing order. Every occupied level \(h>0\) has a predecessor at each
level 1,\ldots,h-1 along a shortest path. Thus \(h_{i}\le i\) for
\(i=1..z\), and the monotonicity of \(U_{h}\) implies the stronger total
bound sum\_\{i=1\}\^{}z \(U_{i}\). The ranked bound is incorporated and
strengthened by density-layer pruning in the final source
\texttt{No89\_v023\_layered.cpp}. The final source also rejects any
component of the current partial graph that contains neither a P vertex
nor an unprocessed off-path vertex. All incident edge choices of such a
component have already been decided, so no connected completion exists.

\subsection{6.6 Row loss, fractional costs and
reflection}\label{row-loss-fractional-costs-and-reflection}

In a partial graph with fixed path and attachments, let
\(D_{xy}=\operatorname{min}(d,dist_{partial}(x,y))\), taking d for
unreachable pairs, and \(w=sum\) D. These are upper bounds on distances
in any admissible completion; the path-path entries equal
\textbar i-j\textbar{} if the path is still geodesic. The current
surplus (not necessarily nonnegative) is r0. Let E? be the undecided
off-path edges still permitted. For \(T=W+(d-1)r\), a universal upper
bound is

\[U_0=w+(d-1)r_0+\sum_{xy\in E_?}(d-D_{xy}).\]

Indeed an optional edge contributes at most d to T after its inclusion
(distance 1 plus d-1), while omission permits at most \(D_{xy}\). Other
distances cannot increase. When deciding a row, deduct \(d-D_{xy}\) for
each omitted row edge. For selected edges xy incident to the row vertex
x, the path-distance caps at x improve to at most min(D\_xp,1+min\_y
D\_yp); at each selected y they improve to at most
\(\operatorname{min}(D_{yp},1+D_{xp})\). Sum these savings over the
distinct path-off pairs, using a minimum cap for x. These pairs are
disjoint from the off-off pairs used in omission penalties. Their total
can be deducted from U0. Discard a branch only if the resulting upper
bound is strictly smaller than the minimum target T over the middle r
range. The final source uses these rules. The earlier row-budget control
is not a separate independent implementation.

At a partial state, let deg?(x) count undecided edges incident to x. For
an optional edge xy define a\_xy=sum\_\{p in P\}
\(\operatorname{max}(0,D_{xp}-1-D_{yp})\). If xy is selected, the total
decrease of path-off distance caps at x in the completion is at least
\(a_{xy}\). For any selected incident set, it is therefore at least sum
\(a_{xy}/deg?(x)\): the maximum of a nonnegative list is at least its
sum divided by the number of all optional incident edges. For each
selected xy the direct pair decrease is \(D_{xy}-1\); the path-off
contributions at its endpoints may be added with these fractions, since
we have bounded the sum separately per x and direct off-off pairs are
disjoint. Thus charge cost

\[c_{xy}=D_{xy}-1+\frac{a_{xy}}{\deg_?(x)}+\frac{a_{yx}}{\deg_?(y)}.\]

For t selected optional edges, w minus the t smallest such costs is an
upper bound on W. Test every attainable middle \(r=r0+t\), retaining
equality. The code scales by 2520 (divisible by every positive degree
\textless=9), and uses integer arithmetic. This rule is used at early
rows only. Reversing the ordered diametral path and relabeling the
off-path vertices preserves all graph data. For each sorted mask tuple,
sort its reflected masks, and skip it if it is lexicographically greater
than its reflection. At least one representative of every orbit remains.
Shape testing explicitly accepts both orientations; changed tie counts
are representation counts, not distinct graph counts.

\subsection{6.7 Exact distances}\label{exact-distances}

For a partial graph containing the path P, if some path distance is
shorter than its index difference, a shortest path realizing such a
shortcut must visit an off-path vertex x. Then for some \(i<j\),
\(\operatorname{dist}(x,p_{i})+\operatorname{dist}(x,p_{j})<j-i\).
Conversely such an inequality supplies a shortcut. With distances capped
at d, this strict test is unchanged: a capped term equal to d cannot
participate in a sum strictly smaller than \(j-i\le d\). Test for each x
in linear time over j, maintaining min\_\{i\textless j\}(D\_xp\_i+i).
Once no shortcut exists, all P-P distances equal their index differences
and contribute binom(d+2,3). BFS from each off-path vertex supplies
every remaining distance, including symmetric path-off entries. At final
exact checks require every off-path BFS to reach every vertex within d;
the path itself already connects its vertices, so this certifies
connectivity and diameter exactly d.~No approximation of the objective
is used at the leaf.

At every surviving edge subset, exact distance checks verify
connectivity, diameter at most d, all P distances equal to coordinate
differences (so diameter exactly d), and compute W with integer
arithmetic. Only candidates capable of reaching the prescribed value are
sent to this final check; inferior graphs are intentionally not counted
as valid. For equality, compare every P-neighbor mask and off-path edge
against the canonical construction, allowing path reversal and
permutations within a group. This is sufficient for isomorphism to the
canonical graph. A noncanonical representation would be reported as an
unresolved tie even if isomorphic by some different map. The final
implementation reports best=pred, positive ties, and noncanonical=0 in
every required row. Tie counts are representation counts, not counts of
distinct graphs.

\subsection{6.8 Density layers}\label{density-layers}

Let z be the number of off-path vertices with empty P-neighbor mask; the
nonempty-mask vertices form the level-zero set of size k-z. In any
connected completion, each empty-mask vertex has finite distance
\(h\ge 1\) to that set in the off-path graph (a path to P first
encounters a nonempty-mask vertex). The levels \(1,...,\ell\) are all
nonempty. An edge joins vertices whose levels differ by at most one, by
the triangle inequality for distance to a set. Therefore for level sizes
s0=k-z,s1,\ldots,sell, sum\_\{i\textgreater=1\}si=z,

\[m_{\mathrm{off}}\le\sum_i\binom{s_i}{2}+\sum_{i=0}^{\ell-1}s_i s_{i+1}.\]

At a fixed mask tuple with a total of a attachments to P, every
admissible middle-range completion requires \(m_{off}\ge C_{k}+1+k-a\),
where \(C_{k}=A-\operatorname{floor}(k^2/4)\). Enumerate every ordered
positive composition of z, retaining exactly those profiles whose
capacity upper bound reaches this lower bound. For \(z=0\) use the empty
profile; \(z=k\) has no connected completion. This is a necessary
condition, not an assertion that a profile is realizable. There are at
most \(2^(z-1)\) positive compositions for \(z\ge 1\).

Here h denotes the distance level, not the surplus parameter. For an
empty-mask vertex at level h, distance to \(p_{j}\) is at most
min(d,h+md(m,j)) for some nonempty mask m. Hence its total path-distance
sum is at most \(U_{h}=max\) over nonempty m of \(sum_{j}\)
min(d,h+md(m,j)). \(U_{h}\) is nondecreasing. For any retained profile
with sorted levels \(h1\le ...\le hz\), the total is bounded by sum
\(U_{hi}\). For a partial graph let the existing distance-sum caps of
the empty vertices, sorted, be \(v1\le ...\le vz\). The stronger bound
is \(sum_{i}\) \(\operatorname{min}(v_{i},U_{hi})\). To see why the
sorted pairing is safe, for \(a\le b\) and \(c\le d\),
\(\operatorname{min}(a,c)+\operatorname{min}(b,d)\ge \operatorname{min}(a,d)+\operatorname{min}(b,c)\).
Repeated swaps show that pairing in sorted order maximizes the sum among
permutations. Thus the unknown association of vertices to levels cannot
exceed this value. Maximize over retained profiles. A profile
componentwise dominated by another retained profile can be removed
because every \(U_{h}\) is nondecreasing and all min expressions are
monotone. At the mask root and at partial row states, replace just the
empty vertices' path-distance contribution by the appropriate upper
bound. Keep w uncorrected in the separate row-loss and fractional edge
bounds, to avoid double counting path-distance savings. No assumption
about unknown edges beyond connectivity, the required minimum edge count
and the distance-to-set triangle inequality is made.
\texttt{No89\_v023\_layered.cpp} implements this exact relaxation. All
prunes retain possible equality.

\section{7. Sharp uniform threshold}\label{sharp-uniform-threshold}

Choose an integer j maximizing j(k-2j) for
\(1\le j\le \operatorname{floor}(k/2)\). Completing the square gives
\(j(k-2j)=k^2/8-2(j-k/4)^2\). Write \(k=4a+b\), b in \{0,1,2,3\}; choose
a nearest integer to k/4. The distance is c/4 with
\(c=\operatorname{min}(b,4-b)\), and \(k^2\) mod \(8=c^2\). Thus the
maximum is \(M=\operatorname{floor}(k^2/8)\); the nearest integer is in
the required range for \(k\ge 4\).

Set \(h*=(j-1)(k-j+1)+1\). It is positive. For \(k\ge 8\),
\(x=j-1\le \operatorname{floor}(k/4)\le \operatorname{floor}(k/2)-2\).
Since x(k-x) increases on this interval,
\(h*\le \operatorname{floor}(k^2/4)-3\). For \(k=4,5,6,7\) the possible
h* are respectively 1, 1, 1 or 6, and 7, all strictly below
\(\operatorname{floor}(k^2/4)\). The minimal feasible R0 is j because
the preceding product is h\emph{-1 and the product at j is h}+k-2j.
Hence \(L_0=k-j\), \(q=k-2j\ge 1\) and \(qR0=M\).

At any \(d\ge 4\), construct a rival V with the same path, left and
right cliques of sizes k-j+1 and j-1, all ordinary masks \{0,1\} in
local orientation, except one left vertex x* whose p0 attachment is
removed. Both graphs have d+1+k vertices. Relative to \(H_*\), V gains
k-2j+1 clique edges and loses q+1 attachments, so its edge count and r
are identical. The path remains geodesic: same-side attachments span at
most two and no cross-side edges exist. All distances are at most d:
left and right ordinary vertices meet via p1 and \(p_(d-1)\) at distance
d; x* also has p1 as a neighbor and distance 2 to p0. Shortened \(H_*\)
vertices only reduce these distances. Both diameters are d.

In V the unique abnormal mask has \(e=0\), \(a=1\), \(\sigma =1\),
\(T=2\), \(U=1\). The two sides are cliques, with \(E=J=0\) and
\(\delta =1\). The identity of Section 3 gives

\[W(V)-W(H_*)=M+2-d.\]

At \(d=M+2\) (which is at least 4), the two graphs tie. In \(H_*\) the
two path endpoints and all k-q ordinary off-path vertices have
eccentricity d; the q shortened vertices have eccentricity d-1. In V all
k off-path vertices and the two endpoints have eccentricity d.~The
internal path vertices in both have eccentricity below d: their
distances to the endpoint groups are at most max(i,d-i), with any
same-end exception at most 2. Thus the counts are k+2-q and k+2, which
are different. Eccentricity is invariant under isomorphism, so V is not
isomorphic to \(H_*\). At this diameter either \(H_*\) is not maximal or
it is not the unique maximizer. The simultaneous value-and-uniqueness
assertion therefore fails at M+2. Sections 2--6 prove it for all
\(d\ge M+3\). This gives the exact uniform threshold. The threshold for
the value alone is not settled by this witness at M+2; we make no claim
about it or about an individual h threshold.

\section{8. Computational evidence and
limitations}\label{computational-evidence-and-limitations}

The final finite certificate is \texttt{No89\_v023\_layered.cpp}
(C++17), unchanged from v0.23.
\texttt{No89\_v019\_integer\_certificate.py} implements Section 5;
\texttt{check.py} checks formula, coverage, exit markers, equality
counts and small lower witnesses, but does not prove C++ pruning safety.
\texttt{run\_all.py} builds and executes the finite certificate. The
branch-bound proof in Section 6 and the independent AI audit address
pruning safety, not the log checker. Historical exploratory sources and
interrupted controls are not required for reproduction.

\begin{longtable}[]{@{}rcr@{}}
\toprule\noalign{}
k & Required diameters & Rows \\
\midrule\noalign{}
\endhead
\bottomrule\noalign{}
\endlastfoot
4 & 5--8 & 12 \\
5 & 6--10 & 25 \\
6 & 7--12 & 48 \\
7 & 9--14 & 66 \\
8 & 11--16 & 90 \\
9 & 13--18 & 114 \\
10 & 15--20 & 144 \\
Total & 39 instances & 499 \\
\end{longtable}

\begin{longtable}[]{@{}rrrr@{}}
\toprule\noalign{}
k & d0 & Relaxed states & Minimum doubled margin \\
\midrule\noalign{}
\endhead
\bottomrule\noalign{}
\endlastfoot
11 & 18 & 48928 & 20 \\
12 & 21 & 74782 & 30 \\
13 & 24 & 104284 & 36 \\
14 & 27 & 152455 & 42 \\
15 & 31 & 207943 & 50 \\
16 & 35 & 293315 & 58 \\
17 & 39 & 391835 & 66 \\
18 & 43 & 539897 & 74 \\
\end{longtable}

All 499 rows were completed on Linux and reproduced on Windows (MinGW
g++ 13.2.0, Python 3.12.10). The independent AI auditor also reported a
matching rerun. The final audit was CONDITIONAL PASS: the analytic proof
and integer certificate were assessed PASS within their respective
domains; the finite computation retained conditions on implementation
evidence. Stage 5 reported 168/168 small differential checks for
\(k=4\), 98/210 for \(k=5\), and 128/216 larger pruning-disable
comparisons; all completed comparisons agreed. Independent geng
enumeration and adversarial search were partial, not complete.
ASan/UBSan checks covered 33/39 diameter instances \((k=4..9)\);
\(k=10\) was not run under sanitizers. Static bounds were assessed, but
do not replace dynamic checks. No counterexample or mathematical FAIL
was reported. This is an AI-assisted, computer-assisted research
preprint, not a human-refereed result or a machine-checked formal proof.
The \(k=10\) integer relaxation has 44 unresolved relaxed states and is
not used as a proof for \(k=10\). Auxiliary tests support, but do not
replace, the mathematical argument. Original Linux logs lack source hash
headers; they are preserved unchanged. Windows logs contain source and
binary hashes; the independent audit reported a rerun. See the
reproduction package for the separate provenance record.

\section{9. Related work and scope of
novelty}\label{related-work-and-scope-of-novelty}

Borovićanin et al.~{[}1, Theorem 2.13{]} identify balanced double brooms
as the unique maximizing trees for sufficiently large order in their
stated diameter range, and Proposition 2.15 supplies triple-broom
competitors. Their theorem concerns trees \((r=0)\); the strict middle
range here requires positive surplus and fixes the edge count, so the
two statements have different domains. The similar quadratic scale alone
is not a novelty claim. The fixed-order and fixed-diameter problem
without a prescribed edge count is also distinct from the present
problem. Sun et al.~{[}2{]} address maximum Wiener index at given
diameter; Cambie {[}3{]} treats an asymptotic problem of Plesník.
References {[}2,3{]} are included as related background identified in
{[}1{]}, not asserted here to have been fully re-audited. The present
paper makes no claim to solve the unrestricted problem, the tree
threshold, or all fixed-edge regimes. The literature comparison is
limited; novelty remains inconclusive and references partial.

\section{References}\label{references}

{[}1{]} B. Borovićanin, D. Božović, E. Glogić, D. Mesarič Štesl, S.
Špacapan and E. Zogić, \emph{New results on Wiener index of trees with a
given diameter}, Quaestiones Mathematicae (2026), 1--16, online 11
September 2026. DOI: 10.2989/16073606.2026.2718457.

{[}2{]} Q. Sun, B. Ikica, R. Škrekovski and V. Vukašinović, \emph{Graphs
with a given diameter that maximise the Wiener index}, Applied
Mathematics and Computation 356 (2019), 438--448.

{[}3{]} S. Cambie, \emph{An asymptotic resolution of a problem of
Plesník}, Journal of Combinatorial Theory, Series B 145 (2020),
341--358.

\textbf{Status:} NOT HUMAN PEER REVIEWED / NOVELTY INCONCLUSIVE /
REFERENCES PARTIAL. Independent AI audit: CONDITIONAL PASS. Copyright
2026 Ryutaro Yonezu. All rights reserved.

\end{document}
